\documentclass{article}
\usepackage[utf8]{inputenc}
\usepackage[T1]{fontenc}
\usepackage{geometry}
\usepackage{graphicx}
\usepackage{float}
\usepackage{hyperref}
\usepackage[english,indonesian]{babel}

\title{Optimasi Pemrosesan Template LaTeX Menggunakan Pendekatan Berbasis AI}
\author{Academic Support Agent}
\date{\today}

\begin{document}
\maketitle

\begin{abstract}

\end{abstract}

\chapter{PENDAHULUAN}


\chapter{TINJAUAN PUSTAKA}


% Original content (commented out):
% \documentclass[12pt, a4paper]{article}
\usepackage[brazilian]{babel}
\usepackage[utf8]{inputenc}
\usepackage{pslatex} %Fonte times new roman
\usepackage{amsmath}
\usepackage{amsfonts}
\usepackage{amssymb}
\usepackage{amsthm}
\usepackage{makeidx}
\usepackage{graphicx}
\usepackage{booktabs}
\usepackage{multicol}
\usepackage{float}
\usepackage{fancyhdr} 
\usepackage{etoolbox}
\pagestyle{fancy} 
\usepackage[portuguese, ruled, linesnumbered]{algorithm2e}
\usepackage[left=2cm,right=2cm,top=2cm,bottom=2cm]{geometry}

\renewcommand{\qedsymbol}{$\blacksquare$}

\newtheorem{definicao}{Definição}
\newtheorem{teorema}{Teorema}
\newtheorem{lema}{Lema}
\newtheorem{proposicao}{Proposição}


%% É importante que não seja inserido novos pacotes neste documento. Ele deve ser feito nesses parâmetros.

\makeatletter
\renewenvironment{thebibliography}[1]
     {\section*{\bibname}% <-- this line was changed from \chapter* to \section*
      \@mkboth{\MakeUppercase\bibname}{\MakeUppercase\bibname}%
      \list{\@biblabel{\@arabic\c@enumiv}}%
           {\settowidth\labelwidth{\@biblabel{#1}}%
            \leftmargin\labelwidth
            \advance\leftmargin\labelsep
            \@openbib@code
            \usecounter{enumiv}%
            \let\p@enumiv\@empty
            \renewcommand\theenumiv{\@arabic\c@enumiv}}%
      \sloppy
      \clubpenalty4000
      \@clubpenalty \clubpenalty
      \widowpenalty4000%
      \sfcode`\.\@m}
     {\def\@noitemerr
       {\@latex@warning{Empty `thebibliography' environment}}%
      \endlist}
\makeatother


\begin{document}
%% Neste modelo, a formatação já está contida como pede para a submissão do trabalho, basta reescrever para o seu trabalho %%
%% Mais informações úteis para a construção e utilização do LaTex para seu trabalho estão dispostas nesta janela de edição de código, esses comentários estão apresentados na seguinte formatação: %% ... %% . Como dito no documento, caso tenha mais dúvidas, entre em contato pelo e-mail smat.fct@unesp.br %%

%%%%%%%%%%% Cabeçalho %%%%%%%%%%%%%%

\addtolength{\headwidth}{90pt}	
	\addtolength{\headsep}{-15pt}
	\renewcommand{\headrulewidth}{0pt}	
	\vspace*{0.1cm}		
	\fancyhead[R]{{\footnotesize Presidente Prudente (SP), de 27 a 30 de agosto de 2024.}}
	\fancyhead[L]{{\footnotesize XVIII SMAT}}
	
	\makeatletter
	\def\footrule{{\if@fancyplain\let\footrulewidth\plainfootrulewidth\fi
			\vskip-\footruleskip\vskip-\footrulewidth
			\leavevmode\rlap{\hspace*{-0.05in}\rule{0.25\paperwidth}{\footrulewidth}}
			\vskip\footruleskip}}

\makeatother
	\fancyfoot{}
	\fancyfoot[R]{}
	\fancyfoot[L]{}

%%%%%%%%%%%%%%%%%%%%%%%%%%%%%%%%%

%%%%%%%%%%%% Título do trabalho e autores %%%%%%%%%%%%%%%%%%%%%%
\begin{center}
	{\large\textbf{INSTRUÇÕES PARA PREPARAÇÃO DE RESUMO EXPANDIDO E \\
	                    APRESENTAÇÃO DO TRABALHO NO XVIII SMAT\footnote{\scriptsize{XVIII Simpósio de Matemática da Faculdade de Ciências e Tecnologia.}}}}
\vspace*{0.5cm}	
	
	\textbf{Alberto D. Caeiro,  Álvaro M. de Campos\footnote{\scriptsize{Bolsista de Iniciação Científica da FAPESP.}}}\\
    {\footnotesize Departamento de Matemática e Computação, FCT, UNESP\\
	19.060-900, Presidente Prudente, SP\\
	\{adc, amc\}@fct.unesp.br}
	
\vspace*{0.5cm}	

	\textbf{Ricardo C. Reis}\\
	{\footnotesize UNOEXTRA -- Faculdade de Ciências \\
	Rua Fernando Pessoa, 1936 \\
	19.188-700, São José, SP \\
	rcr@jes.porto.pt}
\end{center}
\vspace*{-0.5cm}
%%%%%%%%%%%%%%%%%%%%%%%%%%%%%%%%%%%%%%%%%%%%%%%%%%%%%%%%%%%%

%%%%%%%%%%%%%%%%%%%% INICIO DO TRABALHO %%%%%%%%%%%%%%%%%%%%
\noindent \rule{17cm}{0.8pt}

\vspace*{0.3cm}

{\small \textbf{Resumo:} : Este texto descreve os estilos de formatação que devem ser utilizados na preparação dos resumos a serem submetidos ao XVIII Simpósio de Matemática da Faculdade de Ciências e Tecnologia da UNESP , a ser realizado no período de 27 a 30 de agosto de 2024. Os nomes e endereços dos autores destas instruções são fictícios. Caso a formatação do artigo final não atenda às recomendações da Comissão Organizadora, o trabalho poderá ser rejeitado. A Comissão sugere que o(s) autor(es) faça(m) uma cópia deste arquivo e aproveite(m) a formatação existente para a elaboração do trabalho a ser enviado. Todos os trabalhos submetidos que forem aceitos serão apresentados de forma oral e presencial, na forma de pôster digital. O tempo de cada apresentação será da seguinte forma: autores de trabalhos em nível iniciação científica júnior (estudantes do ensino médio) e iniciação científica (onde o autor-apresentador deverá no mínimo estar cursando a graduação) será de 10 minutos; trabalhos em nível pós-graduação (onde o autor-apresentador deverá no mínimo estar cursando o mestrado ou doutorado) será de 15 minutos. Em ambos os casos, ao final, haverá outros 5 minutos abertos a perguntas.
	\vspace*{0.3cm}	
	\par\textbf{Palavras-Chave:} \textit{Norma, Submissão de resumo, Apresentação do trabalho.}}
	

\noindent \rule{17cm}{0.8pt} 
%%%%%%%%%%%%%%%%%%%%%%%%%%%%%%%%%%%%%%%%%%%%%%%%%%%%%%%%%%%%

\begin{multicols}{2}

\section{Introdução}
\noindent O título do trabalho deve ser digitado em letras maiúsculas, fonte Times New Roman, tamanho 14 pt e em negrito com alinhamento centralizado, e o nome do autor, co-autores e orientador (quando for o caso), também em negrito, tamanho 12 pt, seguido do endereço eletrônico, tamanho 8 pt, logo após o título do trabalho.
\\ A identificação do(s) autor(es) é permitida apenas logo após o título. O texto não pode conter qualquer tipo de identificação, ou seja, não é permitida qualquer referência explícita da autoria no corpo ou em rodapé e/ou nome de projeto/pesquisa, grupo de pesquisa etc. Caso ocorra a identificação em forma ou local diferente do autorizado, o trabalho será excluído por quebra de anonimato.
\\ Abaixo da filiação dos autores e antes do início do artigo deverá ser colocado um resumo, utilizando fonte Times New Roman, tamanho 11 pt, espaçamento simples entre linhas e alinhamento justificado, que deve conter entre 100 e 300 palavras e um recuo de 0,5cm, em ambos os lados. Após o resumo, pule uma linha e introduza no mínimo uma e no máximo, três palavras-chave, separadas por vírgula e com a primeira letra de cada uma delas em maiúscula. Observem que, tanto o resumo como as palavras-chave estão limitados superior e inferiormente por uma linha horizontal de 17 cm. 
\\No rodapé da primeira página deve ser informada a agência de fomento, se houver.
O texto deve conter detalhes suficientes para expor o problema, o método ou formulação utilizados e os resultados obtidos. Deve ser escrito em língua portuguesa e conter INTRODUÇÃO, METODOLOGIA, RESULTADO(S), CONCLUSÃO(ÕES) e REFERÊNCIAS BIBLIOGRÁFICAS, em no mínimo 3 páginas e, no máximo 4 páginas.

\section{Metodologia}

\noindent O trabalho deve ser digitado em LaTex no editor de texto on-line Overleaf, utilizando fonte Times New Roman, tamanho 12 pt, apresentado em duas colunas, de mesma dimensão (8,00 cm de largura, com separação de 1,0 cm entre elas), no formato A4, margens superior e inferior de 2,5 cm, margens esquerda e direita de 2,0 cm, espaçamento simples entre linhas e alinhamento justificado. Entre o título e o(s) autor(es), bem como entre este(s) e o corpo do texto, é necessário espaço duplo.
\\ O trabalho deve apresentar resultados de pesquisa e de estudos realizados pelos discentes e docentes de cursos de Licenciatura e/ou Bacharelado em Matemática e por demais pesquisadores que atuam na área de Matemática Pura, Matemática Aplicada, Educação Matemática e suas áreas.

\section{Resultados}
\noindent As referências devem estar ao final do texto, contendo exclusivamente as obras citadas, em ordem alfabética, pelo sobrenome do primeiro autor. Como exemplo, a referência bibliográfica \cite{3} refere-se a livro, a referência  \cite{2} a artigo em revista e a \cite{1} refere-se a artigo publicado em ``proceedings" de congresso, conforme as normas técnicas da ABNT, NBR 6023, de agosto de 2002.
O texto pode conter figuras, equações e tabelas. No caso da figura, tanto a numeração, quanto o título, devem estar dispostos abaixo desta. A Figura \ref{figura1} é um exemplo.

%%A adição de imagens é feita pelo overleaf. Para acrescentar imagem, aperte no botão esquerdo superior do editor, que exibe a mensagem "Carregar" quando o mouse passa por cima. Depois aperte em "Select Files" e selecione os arquivos das imagens no seu computador que serão adicionados no documento%%

\begin{figure}[H]
\centering
\resizebox{\columnwidth}{!}{\includegraphics[scale=0.4]{figura1.pdf}}
%%%%%% Use o comando acima para que as imagens incluidas fiquem com a largura estritamente igual a largura da coluna
%\includegraphics[scale=0.4]{figura1.pdf}
%%%%%% Use este comando ao invés do outro para definir a escala da imagem incluida como quiser.
\caption{Função {\it seno}.} 
\label{figura1}
\end{figure}

%% A imagem deve ser inserida em \includegraphics[]{nome.png}, onde você deverá substituir nome.png pelo nome e extensão respectivos do arquivo que foram adicionados no overleaf%%

\noindent As equações devem estar centralizadas, com as suas respectivas numerações entre parênteses e alinhadas a direita, como mostra o exemplo dado pela equação \eqref{equacao1}. Se a referência à equação ocorrer no início do parágrafo, utilize: ``Equação \eqref{equacao1} mostra a relação ... .''

\begin{equation}\label{equacao1}
     E = mc^2.                    
\end{equation}

\noindent  A Tabela \ref{tabela1} mostra os tamanhos e tipos de fontes utilizados. Observe que a numeração e o nome da tabela estão dispostos acima da tabela. Entre o texto, figuras, equações e tabelas, deixe sempre uma linha em branco.
\begin{table}[H]
\centering
\caption{Tamanho das fontes.}
{\footnotesize
\begin{tabular}{lll}
 \hline
Texto & Fonte & Aparência \\ 
\hline                              
Título & 14pt        & {\bf negrito} \\
Autor(es) & 12pt    & {\bf negrito} \\
Endereço(s) & 8pt              & normal \\
Resumo & 11pt       & normal \\
Texto principal & 12pt       & normal\\
\hline
\end{tabular}}
\label{tabela1}
\end{table}

%%O site tablesgenerator.com oferece uma ferramenta mais simplificada para criar ou colar seus dados ou informações que deseja apresentar em tabela em uma planilha facilmente editável, para então gerar automaticamente a tabela desejada em texto para inserir nesta seção de edição do código para colocar sua tabela na apresentação%%

\section{Conclusões}

\noindent A data final de submissão de trabalhos pode ser consultada no site do evento
(www.fct.unesp.br/extensao/smat/edicao-atual---2024/) em “Datas
Importantes”. Os trabalhos aceitos serão publicados no 
formato online com ISSN 2526-3099.
No caso de eventuais correções sugeridas 
pelos pareceristas, o(s) autor(es) será(ão) 
devidamente notificados e terá(ão) um 
prazo adicional para efetuá-las.
\\Somente serão aceitos trabalhos que atendam aos critérios previamente estabelecidos. 
\\Os arquivos enviados para a submissão do trabalho devem ser o pdf do LaTex e a pasta compactada (zip) do LaTex, ambos nomeados da seguinte forma: sobrenome do autor principal, seguido de XVIII SMAT. Por exemplo, no caso do artigo relativo a essas normas, o nome seria:\linebreak ``Caeiro\_XVIIISMAT.pdf''. A nomenclatura do arquivo do LaTex (.tex) também deve ser nomeado desta forma. Em caso da submissão de mais de um trabalho pelo mesmo autor deve-se acrescentar um número na frente do nome.
\\ No caso de eventuais dúvidas em relação a formatação, entrar em contato com a Comissão Organizadora por meio do endereço eletrônico smat.fct@unesp.br.

%% O nome do PDF e da pasta Compactada do LaTex tem mesmo nome do projeto nomeado no Overleaf. Você pode renomear o nome do projeto já no formato "autor_XVIIISMAT" para ser aplicado em ambos os arquivos que você precisa enviar. Para renomear o nome do projeto no overleaf, vá no nome atual do projeto na barra superior acima do editor, e ao passar o mouse por cima, o símbolo de um lápis irá aparecer com o rótulo "Renomear", apertando nele será possível renomear o projeto e simultaneamente os documentos.%%

%% Para renomear o documento do .tex, basta clicar nos três pontos à direita do arquivo "main.tex", e renomear de acordo com a formatação pedida pelo arquivo.

%%Para obter o ZIP do LaTex, vá aos seus projetos, que aparecem ao logar no overleaf como página inicial ou no editor, ao apertar o botaõ no canto superior esquerdo com o símbolo de uma casa. Lá, você vai até o título do projeto de sua apresentação e à direita, dentre os diversos botões, um deles lhe dará a opção "Baixar Arquivo .zip"%%

\begin{thebibliography}{99}
\bibitem{1} CANÇADO, A. L. Toxocomanias de substituição. In: CONGRESSO FARMACÊUTICO E BIOQUÍMICO PAN-AMERICANO, 3., 1989, São Paulo. {\bf Anais}.  São Paulo: Fed. das Assoc. de Farmacêuticos do Brasil, 1990. p. 259-300.
\bibitem{2} REY, L. Problemas de saúde pública: hidatidose humana. {\bf Rev. Roche}, Rio de Janeiro, v. 18, n. 2, p. 298-301, nov.1969.
\bibitem{3} TELLES, P. C. S. {\bf Tubulações industriais:} materiais, projeto e desenho. 7. ed. Rio de Janeiro: LTC, 1989. 384 p. 
\end{thebibliography}	

%% É importante que as referências Bibliográficas estejam na formatação da ABNT. Ou seja, o título da referência deve estar em negrito e além disso, é possível obter as referências neste formato para qualquer artigo ou revista ou livro encontrado em uma pesquisa no Google Scholar, pelo botão "cite" das suas referências.

%% No espaço acima, das referências bibliográficas, você deve colocar cada referência com um \bibitem{n}, onde n é o número da referência, e então é possível referenciar a bibliografia de número "n" no texto usando o comando \cite{n}. %%

\end{multicols}
\end{document}

\end{document}